\chapter{\eh{satellite-antenna} Wireless Empresarial --- WPA3 y 802.1X}

\begin{concepto}
El Wi-Fi empresarial moderno no es el mismo router de casa con una contraseña
compartida. En entornos corporativos, aeropuertos, universidades y hospitales,
la seguridad inalámbrica se construye sobre tres pilares:
\textbf{autenticación individual}, \textbf{cifrado por sesión} y
\textbf{control de acceso por puerto}.

Este capítulo cubre \textbf{WPA3} (el protocolo de seguridad Wi-Fi más moderno)
y \textbf{802.1X} (el estándar de control de acceso a red por puerto, usado tanto
en Wi-Fi como en switches cableados).
\end{concepto}

%% ── 1.1 WPA3 ────────────────────────────────────────────────────────────────
\section{\eh{locked} WPA3 --- Wi-Fi Protected Access 3}

\begin{concepto}
\textbf{WPA2} (802.11i, 2004) usa \textbf{PSK} (Pre-Shared Key): una sola
contraseña compartida que todos los usuarios conocen. Quien sabe la contraseña
entra. Si esa contraseña se filtra, \textit{toda la red queda comprometida}.
Además, un atacante puede capturar el 4-way handshake y atacarlo offline
con diccionario (Hashcat a millones de intentos por segundo).

\textbf{WPA3} (2018) reemplaza PSK por \textbf{SAE} (Simultaneous Authentication
of Equals), también conocido como ``Dragonfly Handshake''. Cada cliente se
autentica individualmente con un protocolo de prueba de conocimiento cero.

\emoji{light-bulb} \textbf{Analogía:}
\begin{itemize}
    \item \textbf{WPA2-PSK} = llave maestra que todos los empleados tienen.
          Si un empleado la pierde (o la fotografía), cualquiera entra.
    \item \textbf{WPA3-SAE} = tarjeta inteligente individual por empleado.
          Cada tarjeta es única; perder una no compromete las demás.
          Además, la tarjeta no ``revela'' el PIN al sistema --- solo demuestra
          que lo conoce (prueba de conocimiento cero).
\end{itemize}
\end{concepto}

\begin{center}
\renewcommand{\arraystretch}{1.5}
\begin{tabular}{>{\bfseries}p{3.8cm} p{4.5cm} p{4.5cm}}
\toprule
\textbf{Característica} & \textbf{WPA2} & \textbf{WPA3} \\
\midrule
Intercambio de llaves
    & PSK (Pre-Shared Key) --- contraseña compartida
    & SAE / Dragonfly --- Diffie-Hellman con prueba de conocimiento cero \\
Forward Secrecy
    & No --- capturar el handshake permite descifrar sesiones pasadas
    & Sí --- cada sesión genera llaves efímeras únicas \\
Ataques offline
    & Vulnerable --- atacante captura el 4-way handshake y ataca con diccionario
    & Resistente --- SAE no permite ataques offline de diccionario \\
Redes abiertas
    & Sin cifrado (aeropuertos, cafeterías en texto claro)
    & OWE (Opportunistic Wireless Encryption) --- cifrado sin contraseña \\
Modo empresarial
    & WPA2-Enterprise con 802.1X/EAP
    & WPA3-Enterprise con 802.1X/EAP + Suite B (criptografía de 192 bits) \\
Estándar IEEE
    & 802.11i (2004)
    & IEEE 802.11-2020 (enmienda) \\
\bottomrule
\end{tabular}
\end{center}

%% ── 1.1.1 SAE ───────────────────────────────────────────────────────────────
\subsection{\eh{key} SAE --- Simultaneous Authentication of Equals}

\begin{concepto}
\textbf{SAE} (RFC 7664 / IEEE 802.11-2020) es un protocolo de establecimiento
de llaves basado en \textbf{Diffie-Hellman sobre curva elíptica} con una
\textbf{prueba de conocimiento cero} (\textit{zero-knowledge proof}).

\textbf{¿Qué significa prueba de conocimiento cero aquí?}
En WPA2, durante el 4-way handshake, el cliente y el AP derivan una PMK
a partir de la contraseña y la intercambian de forma que un atacante que capture
el tráfico puede verificar contraseñas offline. En SAE, el protocolo demuestra
que \textit{ambas partes conocen la misma contraseña} sin que ninguna la transmita
--- ni en texto claro ni en forma derivada atacable offline.

\vspace{0.5em}
\textbf{Forward Secrecy (Secreto hacia adelante):}
Cada sesión Wi-Fi negocia un par de llaves efímeras (de un solo uso).
Si un atacante graba todo el tráfico cifrado de hoy y mañana roba
la contraseña del AP, \textit{no puede descifrar el tráfico de hoy}
porque las llaves efímeras ya no existen. WPA2 no tiene esta propiedad.
\end{concepto}

\begin{cuidado}
\textbf{KRACK Attack (Key Reinstallation Attack, 2017) --- CVE-2017-13077}

Mathy Vanhoef descubrió que en el 4-way handshake de WPA2, un atacante
Man-in-the-Middle podía forzar la \textbf{reinstalación de llaves ya usadas}
manipulando los mensajes del handshake. Al reinstalar una llave con un
nonce repetido, el cifrado CCMP se debilita drásticamente --- en algunos
casos hasta recuperación de texto claro.

\textbf{Impacto:} afectó prácticamente todo dispositivo Wi-Fi con WPA2 del planeta.

\textbf{WPA3/SAE es inmune por diseño:} SAE no usa un 4-way handshake con
reinstalación de llaves. El protocolo Dragonfly genera llaves únicas por sesión
que no pueden ser reinstaladas, eliminando la clase de vulnerabilidad que
KRACK explotó.

\emoji{light-bulb} Parches de WPA2 existen, pero WPA3 es la solución estructural.
\end{cuidado}

%% ── 1.1.2 OWE ───────────────────────────────────────────────────────────────
\subsection{\eh{globe-with-meridians} OWE --- Opportunistic Wireless Encryption}

\begin{concepto}
\textbf{OWE} (RFC 8110) resuelve un problema clásico: las redes Wi-Fi abiertas
(aeropuertos, cafeterías, hoteles) no tienen contraseña, por lo que el tráfico
viaja en texto claro y cualquier persona con un adaptador en modo monitor
puede \textbf{sniffear pasivamente} todo lo que transmites.

\textbf{OWE} mantiene la experiencia de usuario (no pide contraseña, te conectas
directamente) pero cifra automáticamente el enlace inalámbrico mediante un
intercambio \textbf{Diffie-Hellman} entre cliente y AP:

\begin{enumerate}
    \item Cliente y AP intercambian claves públicas efímeras DH durante la asociación.
    \item Derivan una \textbf{PMK} (Pairwise Master Key) compartida.
    \item El 4-way handshake estándar procede para generar llaves de sesión.
    \item Todo el tráfico queda cifrado con AES-CCM, incluso sin contraseña.
\end{enumerate}

\emoji{warning} OWE protege contra \textbf{sniffing pasivo} (alguien escucha tu
tráfico Wi-Fi). No autentica al AP --- un Evil Twin todavía puede responder.
Para autenticación del AP, se necesita 802.1X o certificados adicionales.
\end{concepto}

%% ── 1.2 802.1X ──────────────────────────────────────────────────────────────
\section{\eh{shield} 802.1X --- Port-Based Network Access Control}

\begin{concepto}
\textbf{802.1X} (IEEE 802.1X-2020) es el estándar de \textbf{control de acceso
a red basado en puerto}. Define tres actores:

\begin{itemize}
    \item \textbf{Supplicant} (Suplicante): el cliente que quiere acceso a la red
          (laptop, teléfono, impresora). Implementado en el sistema operativo
          (wpa\_supplicant en Linux, native supplicant en Windows/macOS).
    \item \textbf{Authenticator} (Autenticador): el switch o Access Point.
          Tiene el puerto en estado \textbf{UNAUTHORIZED} (bloqueado) hasta que
          el servidor RADIUS aprueba al cliente. Solo deja pasar tráfico EAP
          (EAPoL) durante la autenticación.
    \item \textbf{Authentication Server} (Servidor de Autenticación): típicamente
          un servidor \textbf{RADIUS} (FreeRADIUS, Cisco ISE, Microsoft NPS).
          Verifica las credenciales y devuelve Access-Accept o Access-Reject.
\end{itemize}

\emoji{light-bulb} \textbf{Analogía:} el sistema de llave magnética de un hotel.
\begin{itemize}
    \item El \textbf{Supplicant} es el huésped que quiere entrar a su habitación.
    \item El \textbf{Authenticator} es la cerradura electrónica de la puerta.
          No importa que el huésped sepa el número de habitación (contraseña
          compartida): sin la tarjeta válida, la puerta no abre.
    \item El \textbf{RADIUS Server} es el sistema central del hotel que verifica
          si esa tarjeta corresponde a una reserva activa y autoriza el acceso.
\end{itemize}

Si la tarjeta no es válida (credenciales incorrectas, certificado expirado,
dispositivo no administrado), el puerto queda en UNAUTHORIZED y el cliente
solo puede acceder a una VLAN de cuarentena --- si es que eso se configura.
\end{concepto}

%% ── 1.2.1 EAP ───────────────────────────────────────────────────────────────
\subsection{\eh{puzzle-piece} EAP --- Extensible Authentication Protocol}

\textbf{EAP} (RFC 3748) es el protocolo de autenticación que viaja dentro de
802.1X. Es ``extensible'' porque define solo el \textit{marco} del intercambio
de autenticación; el mecanismo real es el \textbf{método EAP} elegido:

\begin{center}
\renewcommand{\arraystretch}{1.5}
\begin{tabular}{>{\bfseries}p{2.8cm} p{3.5cm} p{4.2cm} p{2.5cm}}
\toprule
\textbf{Método} & \textbf{Mecanismo} & \textbf{Caso de uso} & \textbf{Estado} \\
\midrule
EAP-TLS
    & Certificados X.509 mutuos (cliente y servidor)
    & Máxima seguridad; dispositivos administrados con PKI propia
    & \textcolor{green!60!black}{\textbf{Recomendado}} \\
PEAP
    & Contraseña dentro de túnel TLS (solo el servidor tiene certificado)
    & Integración con Active Directory; BYOD con credenciales de dominio
    & Aceptable \\
EAP-TTLS
    & Contraseña dentro de túnel TLS con PAP/CHAP interno
    & Linux, Android, BYOD heterogéneo; no requiere certificado en el cliente
    & Aceptable \\
EAP-MD5
    & Hash MD5 de la contraseña (sin cifrado del canal)
    & Legacy; sin protección contra ataques de diccionario offline
    & \textcolor{red}{\textbf{Obsoleto --- No usar}} \\
EAP-FAST
    & PAC (Protected Access Credential) en lugar de certificado
    & Alternativa a PEAP en entornos Cisco sin PKI completa
    & Uso limitado \\
\bottomrule
\end{tabular}
\end{center}

%% ── 1.2.2 RADIUS ────────────────────────────────────────────────────────────
\subsection{\eh{satellite} RADIUS --- Remote Authentication Dial-In User Service}

\textbf{RADIUS} (RFC 2865/2866) es el protocolo de AAA (Authentication,
Authorization, Accounting) más usado en redes empresariales. Opera sobre UDP:

\begin{lstlisting}
  Puertos RADIUS:
    UDP 1812  — Authentication (sustituye al histórico 1645)
    UDP 1813  — Accounting     (sustituye al histórico 1646)

  Mensajes principales:
    Authenticator --> RADIUS:  Access-Request
                               (contiene: User-Name, NAS-IP, EAP-Message,
                                Message-Authenticator [HMAC-MD5])

    RADIUS --> Authenticator:  Access-Accept
                               (+ atributos: Tunnel-Type=VLAN,
                                Tunnel-Medium-Type=802,
                                Tunnel-Private-Group-ID=<VLAN-ID>)

                               Access-Reject
                               (acceso denegado, puede incluir Reply-Message)

                               Access-Challenge
                               (solicita más información al cliente,
                                p.ej. segundo factor, certificado)

  Seguridad:
    - El secreto compartido (RADIUS shared secret) cifra los atributos
      sensibles (contraseñas) usando MD5.
    - Para producción: usar RADSEC (RADIUS over TLS, RFC 6614)
      en lugar de RADIUS sobre UDP plano.
\end{lstlisting}

%% ── 1.2.3 Flujo completo 802.1X ─────────────────────────────────────────────
\subsection{\eh{right-arrow-curving-up} Flujo Completo 802.1X --- Diagrama de Secuencia}

El siguiente diagrama muestra el flujo completo de autenticación 802.1X
con EAP-TLS entre los tres actores:

\begin{center}
\begin{tikzpicture}[
    msg/.style={-{Stealth}, thick},
    msgr/.style={{Stealth}-, thick},
    lbl/.style={font=\small\ttfamily, text=textlight},
    lbls/.style={font=\scriptsize\itshape, text=gray!60!white}
]

%% ── Lifelines ───────────────────────────────────────────────────────────────
\node[font=\bfseries\small, align=center, text=textlight] (sup) at (0, 0)
    {Supplicant\\{\scriptsize\itshape(Laptop)}};
\node[font=\bfseries\small, align=center, text=textlight] (aut) at (6, 0)
    {Authenticator\\{\scriptsize\itshape(Switch / AP)}};
\node[font=\bfseries\small, align=center, text=textlight] (rad) at (12, 0)
    {RADIUS Server\\{\scriptsize\itshape(FreeRADIUS/ISE)}};

%% Lifeline verticals
\draw[thick, gray!40!white] (0,  -0.4) -- (0,  -13.5);
\draw[thick, gray!40!white] (6,  -0.4) -- (6,  -13.5);
\draw[thick, gray!40!white] (12, -0.4) -- (12, -13.5);

%% Activation boxes
\draw[draw=blue!60!white, fill=blue!15!black] (-0.25,-0.4) rectangle (0.25,-13.5);
\draw[draw=orange!60!white, fill=orange!15!black] (5.75,-0.4) rectangle (6.25,-13.5);
\draw[draw=violet!60!white, fill=violet!20!black] (11.75,-0.4) rectangle (12.25,-13.5);

%% State label: port unauthorized
\node[lbls, align=center, text=orange!70!white] at (6, -1.1) {Puerto:\\UNAUTHORIZED};

%% ── Paso 1: EAPOL-Start ─────────────────────────────────────────────────────
\draw[msg, cyan!70!white] (0.25,-1.6) -- (5.75,-1.6)
    node[midway, above, font=\small, text=cyan!80!white] {\textbf{EAPOL-Start}};
\node[lbls, left] at (0,-1.6) {(1)};

%% ── Paso 2: EAP-Request Identity ────────────────────────────────────────────
\draw[msgr, orange!70!white] (0.25,-2.6) -- (5.75,-2.6)
    node[midway, above, font=\small, text=orange!80!white] {\textbf{EAP-Request / Identity}};
\node[lbls, left] at (0,-2.6) {(2)};

%% ── Paso 3: EAP-Response Identity ───────────────────────────────────────────
\draw[msg, cyan!70!white] (0.25,-3.6) -- (5.75,-3.6)
    node[midway, above, font=\small, text=cyan!80!white] {\textbf{EAP-Response / Identity}};
\node[lbls, left] at (0,-3.6) {(3)};
\node[lbl, below right, font=\scriptsize, text=gray!70!white] at (3,-3.6) {user@empresa.com};

%% ── Paso 4: RADIUS Access-Request ───────────────────────────────────────────
\draw[msg, orange!70!white] (6.25,-4.5) -- (11.75,-4.5)
    node[midway, above, font=\small, text=orange!80!white] {\textbf{RADIUS Access-Request}};
\node[lbls, right] at (12,-4.5) {(4)};
\node[lbl, below right, font=\scriptsize, text=gray!70!white] at (7,-4.5) {EAP-Message encapsulado};

%% ── Paso 5: EAP TLS Tunnel negotiation ──────────────────────────────────────
\draw[msgr, violet!80!white] (6.25,-5.5) -- (11.75,-5.5)
    node[midway, above, font=\small, text=violet!90!white] {\textbf{Access-Challenge (TLS ServerHello)}};
\node[lbls, right] at (12,-5.5) {(5)};

\draw[msgr, orange!70!white] (0.25,-6.3) -- (5.75,-6.3)
    node[midway, above, font=\small, text=orange!80!white] {\textbf{EAP-Request (TLS ServerHello)}};

\draw[msg, cyan!70!white] (0.25,-7.1) -- (5.75,-7.1)
    node[midway, above, font=\small, text=cyan!80!white] {\textbf{EAP-Response (ClientCert + Finished)}};

\draw[msg, orange!70!white] (6.25,-7.9) -- (11.75,-7.9)
    node[midway, above, font=\small, text=orange!80!white] {\textbf{Access-Request (ClientCert)}};

\node[lbls, text=yellow!70!white] at (6,-8.6) {\emoji{locked} Túnel TLS establecido};
\node[lbls] at (6,-9.0) {RADIUS verifica certificado X.509 del cliente};

%% ── Paso 6: RADIUS Access-Accept ────────────────────────────────────────────
\draw[msgr, green!60!white, line width=1.5pt] (6.25,-9.7) -- (11.75,-9.7)
    node[midway, above, font=\small\bfseries, text=green!70!white] {\textbf{RADIUS Access-Accept}};
\node[lbls, right] at (12,-9.7) {(6)};
\node[lbl, below right, font=\scriptsize, text=green!60!white] at (7,-9.7)
    {VLAN=10, Session-Timeout=3600};

%% ── Paso 7: EAP-Success ─────────────────────────────────────────────────────
\draw[msgr, green!60!white, line width=1.5pt] (0.25,-10.6) -- (5.75,-10.6)
    node[midway, above, font=\small\bfseries, text=green!70!white] {\textbf{EAP-Success}};
\node[lbls, left] at (0,-10.6) {(7)};

%% ── Paso 8: Puerto autorizado ───────────────────────────────────────────────
\node[lbls, align=center, text=green!60!white] at (6,-11.4) {Puerto:\\AUTHORIZED};
\node[font=\small, text=green!60!white, align=center] at (6,-12.3)
    {\textbf{Cliente asignado a VLAN 10}\\
     \textbf{Tráfico de datos permitido}};

%% Caption note
\node[font=\scriptsize\itshape, text=gray!60!white, align=center] at (6,-13.3)
    {EAPoL (EAP over LAN) entre Supplicant y Authenticator |
     RADIUS/UDP entre Authenticator y Server};

\end{tikzpicture}
\end{center}

\begin{moderno}
\textbf{Google BeyondCorp --- El dispositivo como identidad}

La implementación de 802.1X en Google va más allá del usuario. En el modelo
\textbf{BeyondCorp} (el origen del término ``Zero Trust''), la autenticación
de red usa \textbf{EAP-TLS donde el certificado pertenece al DISPOSITIVO},
no al usuario:

\begin{itemize}
    \item Cada laptop, teléfono y servidor administrado por Google recibe un
          \textbf{certificado X.509 único} emitido por la PKI interna de Google.
    \item El 802.1X valida ese certificado de dispositivo, no una contraseña de usuario.
    \item Si un empleado trae su laptop personal (\textit{dispositivo no administrado}),
          el AP la rechaza en el 802.1X --- no hay RADIUS Access-Accept ---
          independientemente de que el empleado tenga credenciales válidas de usuario.
    \item \textbf{Credenciales de usuario robadas + dispositivo no administrado
          = cero acceso a la red interna}.
\end{itemize}

\emoji{brain} Este modelo invierte el supuesto de confianza tradicional:
en vez de ``confío en quien está dentro de la red'', el modelo es
``\textit{no confío en nadie hasta verificar dispositivo, usuario y contexto}''.
La red es ahora solo un canal de transporte --- el acceso real lo controla
la capa de aplicación con certificados y políticas de dispositivo.
\end{moderno}
